\section{空间运动刚体元件的动能}

\subsection{空间运动刚体元件动能的微元求和形式}

空间运动刚体元件的动能可用微元求和的形式表示为
\begin{equation}
T = \sum_{k=1}^{N} \frac{1}{2}m_k\,\dot{\boldsymbol{r}}_{Ok}^{\mathrm{T}}\dot{\boldsymbol{r}}_{Ok}
\label{eq:8.1}
\end{equation}
式中：
\begin{itemize}
\item $k$ —— 刚体上微元的标号；
\item $N$ —— 微元的总数；
\item $m_k$ —— 微元$k$的质量；
\item $\boldsymbol{r}_{Ok}$ —— 微元$k$在全局惯性坐标系$Oxyz$中的坐标列阵；
\item $\dot{\boldsymbol{r}}_{Ok}$ —— $\boldsymbol{r}_{Ok}$对时间的一阶导数，即在$Oxyz$中的绝对速度列阵。
\end{itemize}

为便于将刚体元件的动能表示为系统广义坐标$\boldsymbol{q}$和广义速度$\dot{\boldsymbol{q}}$的形式，为刚体元件引入一个连体坐标系$R\xi\eta\zeta$，此时$\boldsymbol{r}_{Ok}$可表示为
\begin{equation}
\boldsymbol{r}_{Ok} = \boldsymbol{r}_{OR} + A_{OR}\boldsymbol{l}_{Rk}
\label{eq:8.2}
\end{equation}
式中：
\begin{itemize}
\item $\boldsymbol{r}_{OR}$ —— $R\xi\eta\zeta$连体系原点在全局惯性坐标系中的坐标列阵；
\item $A_{OR}$ —— 从$R\xi\eta\zeta$连体系到全局惯性坐标系的坐标变换矩阵；
\item $\boldsymbol{l}_{Rk}$ —— 微元$k$在$R\xi\eta\zeta$连体系中的坐标列阵。
\end{itemize}

式~(\ref{eq:8.2})对时间的一阶导数可表示为
\begin{equation}
\dot{\boldsymbol{r}}_{Ok} = \dot{\boldsymbol{r}}_{OR} + A_{OR}\tilde{\boldsymbol{\omega}}_{OR|R}\,\boldsymbol{l}_{Rk}
\label{eq:8.3}
\end{equation}
式中$\boldsymbol{\omega}_{OR|R}$表示连体坐标系$R\xi\eta\zeta$相对于全局惯性坐标$Oxyz$的绝对角速度在$R\xi\eta\zeta$系中的投影列阵。

将速度列阵进一步改写为
\begin{align}
\dot{\boldsymbol{r}}_{Ok} &= A_{OR}A_{OR}^{\mathrm{T}}\dot{\boldsymbol{r}}_{OR} + A_{OR}\tilde{\boldsymbol{l}}_{Rk}^{\mathrm{T}}\boldsymbol{\omega}_{OR|R} \nonumber\\
&= \begin{bmatrix}
A_{OR} & A_{OR}\tilde{\boldsymbol{l}}_{Rk}^{\mathrm{T}}
\end{bmatrix}
\begin{bmatrix}
A_{OR}^{\mathrm{T}}\dot{\boldsymbol{r}}_{OR} \\[2pt] \boldsymbol{\omega}_{OR|R}
\end{bmatrix}
=: C_{Rk}\boldsymbol{v}_R
\label{eq:8.4}
\end{align}

\subsection*{用连体坐标系运动学量表示的刚体元件的动能}

将式~(\ref{eq:8.4})代入式~(\ref{eq:8.1})，可得空间运动刚体元件的动能
\begin{align}
T &= \sum_{k=1}^{N} \frac{1}{2}m_k\,\dot{\boldsymbol{r}}_{Ok}^{\mathrm{T}}\dot{\boldsymbol{r}}_{Ok} \nonumber\\
&= \sum_{k=1}^{N} \frac{1}{2}m_k\,(C_{Rk}\boldsymbol{v}_R)^{\mathrm{T}}(C_{Rk}\boldsymbol{v}_R) \nonumber\\
&= \sum_{k=1}^{N} \frac{1}{2}m_k\,\boldsymbol{v}_R^{\mathrm{T}}C_{Rk}^{\mathrm{T}}C_{Rk}\boldsymbol{v}_R \nonumber\\
&= \frac{1}{2}\boldsymbol{v}_R^{\mathrm{T}}\left(\sum_{k=1}^{N} m_k\,C_{Rk}^{\mathrm{T}}C_{Rk}\right)\boldsymbol{v}_R \nonumber\\
&=: \frac{1}{2}\boldsymbol{v}_R^{\mathrm{T}}M_R\boldsymbol{v}_R
\label{eq:8.5}
\end{align}
式中$M_R$称为刚体元件的质量矩阵。

\subsection*{空间运动刚体元件的质量矩阵}

\begin{align}
M_R &= \sum_{k=1}^{N} m_k\,C_{Rk}^{\mathrm{T}}C_{Rk} \nonumber\\
&= \sum_{k=1}^{N} m_k
\begin{bmatrix}
A_{OR} & A_{OR}\tilde{\boldsymbol{l}}_{Rk}^{\mathrm{T}}
\end{bmatrix}^{\mathrm{T}}
\begin{bmatrix}
A_{OR} & A_{OR}\tilde{\boldsymbol{l}}_{Rk}^{\mathrm{T}}
\end{bmatrix} \nonumber\\
&= \sum_{k=1}^{N} m_k
\begin{bmatrix}
A_{OR}^{\mathrm{T}} \\[2pt] \tilde{\boldsymbol{l}}_{Rk}A_{OR}^{\mathrm{T}}
\end{bmatrix}
\begin{bmatrix}
A_{OR} & A_{OR}\tilde{\boldsymbol{l}}_{Rk}^{\mathrm{T}}
\end{bmatrix} \nonumber\\
&= \sum_{k=1}^{N} m_k
\begin{bmatrix}
I_3          & \tilde{\boldsymbol{l}}_{Rk}^{\mathrm{T}} \\[4pt]
\tilde{\boldsymbol{l}}_{Rk} & \tilde{\boldsymbol{l}}_{Rk}\tilde{\boldsymbol{l}}_{Rk}^{\mathrm{T}}
\end{bmatrix}
\label{eq:8.6}
\end{align}

由此可知，此处定义的空间运动刚体元件的质量矩阵$M_R$为定常（时不变）矩阵。

\subsection*{空间运动刚体元件时不变质量矩阵}

\begin{align}
M_R &= \sum_{k=1}^{N} m_k
\begin{bmatrix}
I_3          & \tilde{\boldsymbol{l}}_{Rk}^{\mathrm{T}} \\[4pt]
\tilde{\boldsymbol{l}}_{Rk} & \tilde{\boldsymbol{l}}_{Rk}\tilde{\boldsymbol{l}}_{Rk}^{\mathrm{T}}
\end{bmatrix} \nonumber\\
&= \begin{bmatrix}
I_3\displaystyle\sum_{k=1}^{N}m_k & \displaystyle\sum_{k=1}^{N}m_k\tilde{\boldsymbol{l}}_{Rk}^{\mathrm{T}} \\[14pt]
\displaystyle\sum_{k=1}^{N}m_k\tilde{\boldsymbol{l}}_{Rk} & \displaystyle\sum_{k=1}^{N}m_k\tilde{\boldsymbol{l}}_{Rk}\tilde{\boldsymbol{l}}_{Rk}^{\mathrm{T}}
\end{bmatrix} \nonumber\\
&=
\begin{bmatrix}
I_3\displaystyle\sum_{k=1}^{N}m_k & \left[\displaystyle\sum_{k=1}^{N}m_k\boldsymbol{l}_{Rk}\right]^{\mathrm{T}}_{\times} \\[14pt]
\left[\displaystyle\sum_{k=1}^{N}m_k\boldsymbol{l}_{Rk}\right]_{\times} & \displaystyle\sum_{k=1}^{N}m_k\tilde{\boldsymbol{l}}_{Rk}\tilde{\boldsymbol{l}}_{Rk}^{\mathrm{T}}
\end{bmatrix} \nonumber\\
&=:
\begin{bmatrix}
m I_3              & m\tilde{\boldsymbol{l}}_{RC}^{\mathrm{T}} \\[4pt]
m\tilde{\boldsymbol{l}}_{RC} & \boldsymbol{J}_{R|R}
\end{bmatrix}
\label{eq:8.7}
\end{align}
式中：
\begin{itemize}
\item $m$ —— 刚体元件的质量；
\item $\boldsymbol{l}_{RC}$ —— 刚体元件质心在连体坐标系$R\xi\eta\zeta$中的坐标列阵；
\item $\boldsymbol{J}_{R|R}$ —— 刚体元件在$R\xi\eta\zeta$系中的转动惯量矩阵。
\end{itemize}

\subsection*{矩阵$\boldsymbol{J}_{R|R}$的分量表达式}

刚体元件质量矩阵中的$\boldsymbol{J}_{R|R}$的分量表达式为
\begin{align}
\boldsymbol{J}_{R|R} &= \sum_{k=1}^{N} m_k\,\tilde{\boldsymbol{l}}_{Rk}\tilde{\boldsymbol{l}}_{Rk}^{\mathrm{T}}
 = \sum_{k=1}^{N} m_k\left(l_{Rk}^2 I_3 - \boldsymbol{l}_{Rk}\boldsymbol{l}_{Rk}^{\mathrm{T}}\right) \nonumber\\
&= \begin{bmatrix}
\displaystyle\sum_{k=1}^{N} m_k(\eta_{Rk}^2 + \zeta_{Rk}^2)  & -\displaystyle\sum_{k=1}^{N} m_k\,\xi_{Rk}\eta_{Rk} & -\displaystyle\sum_{k=1}^{N} m_k\,\xi_{Rk}\zeta_{Rk} \\[12pt]
-\displaystyle\sum_{k=1}^{N} m_k\,\eta_{Rk}\xi_{Rk} & \displaystyle\sum_{k=1}^{N} m_k(\xi_{Rk}^2 + \zeta_{Rk}^2)  & -\displaystyle\sum_{k=1}^{N} m_k\,\eta_{Rk}\zeta_{Rk} \\[12pt]
-\displaystyle\sum_{k=1}^{N} m_k\,\zeta_{Rk}\xi_{Rk} & -\displaystyle\sum_{k=1}^{N} m_k\,\zeta_{Rk}\eta_{Rk} & \displaystyle\sum_{k=1}^{N} m_k(\xi_{Rk}^2 + \eta_{Rk}^2)
\end{bmatrix}
\label{eq:8.8}
\end{align}

刚体元件惯性积矩阵：
\begin{equation}
\boldsymbol{I}_{R|R} =
\begin{bmatrix}
\displaystyle\sum_{k=1}^{N} m_k(\eta_{Rk}^2 + \zeta_{Rk}^2)  & \displaystyle\sum_{k=1}^{N} m_k\,\xi_{Rk}\eta_{Rk} & \displaystyle\sum_{k=1}^{N} m_k\,\xi_{Rk}\zeta_{Rk} \\[12pt]
\displaystyle\sum_{k=1}^{N} m_k\,\eta_{Rk}\xi_{Rk} & \displaystyle\sum_{k=1}^{N} m_k(\xi_{Rk}^2 + \zeta_{Rk}^2)  & \displaystyle\sum_{k=1}^{N} m_k\,\eta_{Rk}\zeta_{Rk} \\[12pt]
\displaystyle\sum_{k=1}^{N} m_k\,\zeta_{Rk}\xi_{Rk} & \displaystyle\sum_{k=1}^{N} m_k\,\zeta_{Rk}\eta_{Rk} & \displaystyle\sum_{k=1}^{N} m_k(\xi_{Rk}^2 + \eta_{Rk}^2)
\end{bmatrix}
\label{eq:8.9}
\end{equation}

\section{空间运动刚体元件的重力势能}

空间运动刚体元件的重力势能可用微元求和的形式表示为
\begin{align}
V_{\mathrm{gravity}} &= -\sum_{k=1}^{N} m_k\,\boldsymbol{g}^{\mathrm{T}}\boldsymbol{r}_{Ok} \nonumber\\
&= -\boldsymbol{g}^{\mathrm{T}}\sum_{k=1}^{N} m_k\boldsymbol{r}_{Ok} \nonumber\\
&= -\boldsymbol{g}^{\mathrm{T}}\sum_{k=1}^{N} m_k\left(\boldsymbol{r}_{OR} + A_{OR}\boldsymbol{l}_{Rk}\right) \nonumber\\
&= -\boldsymbol{g}^{\mathrm{T}}\left(m\boldsymbol{r}_{OR} + mA_{OR}\boldsymbol{l}_{RC}\right) \nonumber\\
&= -m\boldsymbol{g}^{\mathrm{T}}\left(\boldsymbol{r}_{OR} + A_{OR}\boldsymbol{l}_{RC}\right)
\label{eq:8.10}
\end{align}

\section{空间运动刚体元件的转动惯量张量}

绕连体坐标系$R\xi\eta\zeta$原点$R$作定点转动的刚体元件的动能可表示为
\begin{align}
T &= \sum_{k=1}^{N} \frac{1}{2}m_k(\vec{\omega}_{OR} \times \vec{r}_{Rk})\cdot(\vec{\omega}_{OR} \times \vec{r}_{Rk}) \nonumber\\
&= \sum_{k=1}^{N} \frac{1}{2}m_k[\vec{r}_{Rk} \times (\vec{\omega}_{OR} \times \vec{r}_{Rk})]\cdot\vec{\omega}_{OR} \nonumber\\
&= \sum_{k=1}^{N} \frac{1}{2}m_k[(\vec{r}_{Rk}\cdot\vec{r}_{Rk})\vec{\omega}_{OR} - (\vec{r}_{Rk}\cdot\vec{\omega}_{OR})\vec{r}_{Rk}]\cdot\vec{\omega}_{OR} \nonumber\\
&= \sum_{k=1}^{N} \frac{1}{2}m_k[\vec{\omega}_{OR}(\vec{r}_{Rk}\cdot\vec{r}_{Rk}) - (\vec{\omega}_{OR}\cdot\vec{r}_{Rk})\vec{r}_{Rk}]\cdot\vec{\omega}_{OR} \nonumber\\
&= \sum_{k=1}^{N} \frac{1}{2}m_k[\vec{\omega}_{OR}\cdot\vec{\vec{I}}(\vec{r}_{Rk}\cdot\vec{r}_{Rk}) - \vec{\omega}_{OR}\cdot\vec{r}_{Rk}\vec{r}_{Rk}]\cdot\vec{\omega}_{OR} \nonumber\\
&= \vec{\omega}_{OR}\cdot\frac{1}{2}\sum_{k=1}^{N} m_k[\vec{\vec{I}}(\vec{r}_{Rk}\cdot\vec{r}_{Rk}) - \vec{r}_{Rk}\vec{r}_{Rk}]\cdot\vec{\omega}_{OR} \nonumber\\
&=: \frac{1}{2}\vec{\omega}_{OR}\cdot\vec{\vec{J}}_R\cdot\vec{\omega}_{OR}
\label{eq:8.11}
\end{align}
式中$\vec{\vec{J}}_R$称为刚体元件相对于点$R$的转动惯量张量。

\subsection*{空间运动刚体元件的转动惯量张量的坐标表示}

刚体元件相对于点$R$的转动惯量张量的坐标表示为
\begin{align}
\vec{\vec{J}}_R &= \sum_{k=1}^{N} m_k[\vec{\vec{I}}(\vec{r}_{Rk}\cdot\vec{r}_{Rk}) - \vec{r}_{Rk}\vec{r}_{Rk}] \nonumber\\
&= \vec{e}_O^{\mathrm{T}}\sum_{k=1}^{N} m_k[I_3(\boldsymbol{r}_{Rk}^{\mathrm{T}}\boldsymbol{r}_{Rk}) - \boldsymbol{r}_{Rk}\boldsymbol{r}_{Rk}^{\mathrm{T}}]\vec{e}_O
   =: \vec{e}_O^{\mathrm{T}}\boldsymbol{J}_{R|O}\,\vec{e}_O \nonumber\\
&= \vec{e}_R^{\mathrm{T}}\sum_{k=1}^{N} m_k[I_3(\boldsymbol{l}_{Rk}^{\mathrm{T}}\boldsymbol{l}_{Rk}) - \boldsymbol{l}_{Rk}\boldsymbol{l}_{Rk}^{\mathrm{T}}]\vec{e}_R
   =: \vec{e}_R^{\mathrm{T}}\boldsymbol{J}_{R|R}\,\vec{e}_R
\label{eq:8.12}
\end{align}
式中$\boldsymbol{J}_{R|O}$与$\boldsymbol{J}_{R|R}$均为实对称矩阵，且满足
\begin{equation}
\boldsymbol{J}_{R|O} = A_{OR}\boldsymbol{J}_{R|R}A_{OR}^{\mathrm{T}}, \qquad
\boldsymbol{J}_{R|R} = \sum_{k=1}^{N} m_k[I_3(\boldsymbol{l}_{Rk}^{\mathrm{T}}\boldsymbol{l}_{Rk}) - \boldsymbol{l}_{Rk}\boldsymbol{l}_{Rk}^{\mathrm{T}}]
= \sum_{k=1}^{N} m_k\,\tilde{\boldsymbol{l}}_{Rk}\tilde{\boldsymbol{l}}_{Rk}^{\mathrm{T}}
\label{eq:8.13}
\end{equation}

\subsection*{空间运动刚体元件的转动惯量张量的平行移轴定理}

刚体元件相对于点$R$的转动惯量张量可借助质心点$C$作如下变换：
\begin{align}
\vec{\vec{J}}_R &= \sum_{k=1}^{N} m_k[\vec{\vec{I}}(\vec{r}_{Rk}\cdot\vec{r}_{Rk}) - \vec{r}_{Rk}\vec{r}_{Rk}] \nonumber\\
&= \sum_{k=1}^{N} m_k\vec{\vec{I}}[(\vec{r}_{RC} + \vec{r}_{Ck})\cdot(\vec{r}_{RC} + \vec{r}_{Ck})]
   - \sum_{k=1}^{N} m_k[(\vec{r}_{RC} + \vec{r}_{Ck})(\vec{r}_{RC} + \vec{r}_{Ck})] \nonumber\\
&= m\vec{\vec{I}}(\vec{r}_{RC}\cdot\vec{r}_{RC}) + \sum_{k=1}^{N} m_k\vec{\vec{I}}(\vec{r}_{Ck}\cdot\vec{r}_{Ck})
   - m(\vec{r}_{RC}\vec{r}_{RC}) - \sum_{k=1}^{N} m_k(\vec{r}_{Ck}\vec{r}_{Ck}) \nonumber\\
&= m[\vec{\vec{I}}(\vec{r}_{RC}\cdot\vec{r}_{RC}) - (\vec{r}_{RC}\vec{r}_{RC})] + \vec{\vec{J}}_C
\label{eq:8.14}
\end{align}
上式即为空间运动刚体元件的转动惯量张量平行移轴定理。

\subsection*{空间运动刚体元件的惯性主轴系}

刚体元件相对于任意点$R$的转动惯量张量在选定的坐标系（如$R\xi\eta\zeta$系）中的投影矩阵为
\begin{align}
\boldsymbol{J}_{R|R} &= \sum_{k=1}^{N} m_k\,\tilde{\boldsymbol{l}}_{Rk}\tilde{\boldsymbol{l}}_{Rk}^{\mathrm{T}}
 = \sum_{k=1}^{N} m_k\left(l_{Rk}^2 I_3 - \boldsymbol{l}_{Rk}\boldsymbol{l}_{Rk}^{\mathrm{T}}\right) \nonumber\\
&= \begin{bmatrix}
\displaystyle\sum_{k=1}^{N} m_k(\eta_{Rk}^2 + \zeta_{Rk}^2)  & -\displaystyle\sum_{k=1}^{N} m_k\,\xi_{Rk}\eta_{Rk} & -\displaystyle\sum_{k=1}^{N} m_k\,\xi_{Rk}\zeta_{Rk} \\[12pt]
-\displaystyle\sum_{k=1}^{N} m_k\,\eta_{Rk}\xi_{Rk} & \displaystyle\sum_{k=1}^{N} m_k(\xi_{Rk}^2 + \zeta_{Rk}^2)  & -\displaystyle\sum_{k=1}^{N} m_k\,\eta_{Rk}\zeta_{Rk} \\[12pt]
-\displaystyle\sum_{k=1}^{N} m_k\,\zeta_{Rk}\xi_{Rk} & -\displaystyle\sum_{k=1}^{N} m_k\,\zeta_{Rk}\eta_{Rk} & \displaystyle\sum_{k=1}^{N} m_k(\xi_{Rk}^2 + \eta_{Rk}^2)
\end{bmatrix} \nonumber\\
&=:
\begin{bmatrix}
J_{\xi\xi} & J_{\xi\eta} & J_{\xi\zeta} \\[2pt]
J_{\eta\xi} & J_{\eta\eta} & J_{\eta\zeta} \\[2pt]
J_{\zeta\xi} & J_{\zeta\eta} & J_{\zeta\zeta}
\end{bmatrix}_{R|R}
\label{eq:8.15}
\end{align}

一般情况下$\boldsymbol{J}_{R|R}$为非对角矩阵。我们总是可以找到这样一个坐标系$R'$，使得$\vec{\vec{J}}_R$在该系中的投影矩阵$\boldsymbol{J}_{R|R'}$为对角矩阵。

对于$3\times 3$的实对称矩阵$\boldsymbol{J}_{R|R}$，由矩阵特征值相关理论可知，一定存在三个线性独立且模为1的$3\times 1$的实数列阵$\boldsymbol{\varphi}_i$，使得
\begin{equation}
\boldsymbol{J}_{R|R}\,\boldsymbol{\varphi}_i = \lambda_i\boldsymbol{\varphi}_i
\label{eq:8.16}
\end{equation}
式中实数$\lambda_i$称为矩阵$\boldsymbol{J}_{R|R}$的特征值，特征向量$\boldsymbol{\varphi}_i$满足两两正交：
\begin{equation}
\boldsymbol{\varphi}_i^{\mathrm{T}}\boldsymbol{\varphi}_j = \delta_{ij}, \quad
\boldsymbol{\varphi}_3 = \boldsymbol{\varphi}_1 \times \boldsymbol{\varphi}_2
\end{equation}

由式~(\ref{eq:8.16})可得
\begin{equation}
\boldsymbol{J}_{R|R}
\begin{bmatrix}
\boldsymbol{\varphi}_1 & \boldsymbol{\varphi}_2 & \boldsymbol{\varphi}_3
\end{bmatrix}
=
\begin{bmatrix}
\boldsymbol{\varphi}_1 & \boldsymbol{\varphi}_2 & \boldsymbol{\varphi}_3
\end{bmatrix}
\begin{bmatrix}
\lambda_1 & 0      & 0      \\
0      & \lambda_2 & 0      \\
0      & 0      & \lambda_3
\end{bmatrix}
\label{eq:8.17}
\end{equation}

上式可简记为
\begin{equation}
\boldsymbol{J}_{R|R}\Phi = \Phi\,\mathrm{diag}(\boldsymbol{\lambda})
\end{equation}
式中
\begin{equation}
\Phi = \begin{bmatrix}
\boldsymbol{\varphi}_1 & \boldsymbol{\varphi}_2 & \boldsymbol{\varphi}_3
\end{bmatrix}, \quad
\Phi^{-1} = \Phi^{\mathrm{T}}, \quad
\boldsymbol{\lambda} = \begin{bmatrix}
\lambda_1 \\ \lambda_2 \\ \lambda_3
\end{bmatrix}
\end{equation}

由$\boldsymbol{J}_{R|R}\Phi = \Phi\,\mathrm{diag}(\boldsymbol{\lambda})$和$\Phi^{-1} = \Phi^{\mathrm{T}}$可得
\begin{equation}
\mathrm{diag}(\boldsymbol{\lambda}) = \Phi^{\mathrm{T}}\boldsymbol{J}_{R|R}\Phi
\label{eq:8.18}
\end{equation}

上式可改写为
\begin{equation}
\boldsymbol{J}_{R|R'} = A_{RR'}^{\mathrm{T}}\boldsymbol{J}_{R|R}A_{RR'}
\label{eq:8.19}
\end{equation}
式中$\boldsymbol{J}_{R|R'}$为对角矩阵，$A_{RR'}$为从$R'$系到$R$系的坐标变换矩阵，$R'$系称为转动惯量张量$\vec{\vec{J}}_R$的惯性主轴系。
